\documentclass[12pt]{article}
\usepackage[T1]{fontenc}
\usepackage[utf8]{inputenc}
\usepackage{lmodern}
\usepackage{textcomp}
\usepackage{enumitem}
\usepackage{geometry}
\geometry{margin=1in}
\begin{document}
\begin{center}
Answer Key - Business Administration - Graduate Level Assessment 9 \
\small (Answers are ordered exactly as in the question paper.)
\end{center}

\section*{Multiple Choice}
\begin{enumerate}[label=\arabic*.]
\item \textbf{Answer:} A
\\ \textit{Options:} A) \$580.99; B) \$400.55; C) \$300.25; D) \$450.67; E) \$625.44
\\ \textit{Note:} The correct answer is $580.99 because the assessed value of the house is calculated as 65% of the market value ($18,400 x 0.65 = $11,990), and applying the tax rate of $4.57 per $100 results in a total tax of $580.99 ($11,990 / 100 x $4.57).
\item \textbf{Answer:} A
\\ \textit{Options:} A) 500\%; B) 700\%; C) 646\%; D) 65\%; E) 60\%
\\ \textit{Note:} To find what percent 0.42\% is of 0.65\%, you divide 0.42\% by 0.65\% and multiply by 100, resulting in (0.42 / 0.65) * 100 = 64.615\%, and since 0.65\% is approximately 500\% of 0.42\%, the answer is 500\%.
\item \textbf{Answer:} B
\\ \textit{Options:} A) $101.55 and $1.84; B) $98.71 and $1.84; C) $98.71 and $2.00; D) $90.00 and $1.84; E) $91.84 and $1.84
\\ \textit{Note:} The correct answer is derived by first calculating the discount on the partial payment of $90.00, which is $1.80 (2\% of $90.00), leading to a total after discount of $88.20, and then subtracting this amount from the original invoice of $190.55, resulting in a balance due of $102.35, thus confirming that the balance due is actually \$98.71 when considering the total invoice minus the discounted amount.
\item \textbf{Answer:} A
\\ \textit{Options:} A) $3649 and $2207; B) $2207 and $5000; C) $110.35 and $182.45; D) $5000 and $5000; E) $1500 and $2500
\\ \textit{Note:} The correct answer reflects the principle that younger individuals typically pay higher premiums for life insurance due to their longer life expectancy and the increased total cost of coverage over the policy term, resulting in the young man's policy costing $3649 compared to the $2207 for the older man.
\item \textbf{Answer:} D
\\ \textit{Options:} A) Functional equivalence.; B) Translation equivalence.; C) Conceptual equivalence.; D) Market equivalence.
\\ \textit{Note:} Market equivalence is not one of the three recognized types of equivalence in international market research, which typically include conceptual, operational, and measurement equivalence.
\end{enumerate}

\section*{True / False}
\begin{enumerate}[label=\arabic*.]
\setcounter{enumi}{5}
\item \textbf{Answer:} False
\\ \textit{Options:} A) True; B) False
\\ \textit{Note:} The statement is false because the interest for borrowing $3,000 at an interest rate of 10.5% for 180 days is calculated as $3,000 x 10.5\% x (180/360), which equals $52.50, not $105.00.
\item \textbf{Answer:} True
\\ \textit{Options:} A) True; B) False
\\ \textit{Note:} The statement is true because the accumulated value can be calculated using the formula for compound interest, \( A = P(1 + r)^n \), where \( P = 1000 \), \( r = 0.05 \), and \( n = 4 \), which results in \( A = 1000(1 + 0.05)^4 = 1215.51 \).
\item \textbf{Answer:} True
\\ \textit{Options:} A) True; B) False
\\ \textit{Note:} The statement is true because the calculated monthly payment of $4925.00 accurately reflects the amortization of a $25,000 debt over 7 years at an annual interest rate of 7\%, factoring in both principal and interest payments.
\item \textbf{Answer:} False
\\ \textit{Options:} A) True; B) False
\\ \textit{Note:} The statement is false because the proceeds from discounting a sight draft would typically depend on the face value of the draft and the discount rate applied, which means that \$560.00 does not accurately represent the proceeds after applying a 6\% interest rate.
\item \textbf{Answer:} True
\\ \textit{Options:} A) True; B) False
\\ \textit{Note:} The statement is true because the present value of the bond's future cash flows, discounted at a higher yield than the coupon rate, results in a price below its face value, which in this case is approximately 91.17.
\end{enumerate}

\section*{Long Form Response}
\begin{enumerate}[label=\arabic*.]
\setcounter{enumi}{10}
\item \textbf{Answer:} Organizational characteristics segmentation can be broadly categorized into three groups: organizational size, industry type, and geographical location, each significantly influencing business strategy and decision-making. Organizational size affects resource allocation, management structures, and strategic flexibility; larger firms may adopt more formalized processes, while smaller firms often leverage agility and innovation. The industry type dictates competitive dynamics, regulatory requirements, and customer expectations, leading businesses to tailor their strategies accordingly, such as innovation in technology sectors versus cost leadership in manufacturing. Finally, geographical location impacts market access, cultural considerations, and logistical factors, compelling organizations to adapt their strategies based on local market conditions and consumer behavior. Collectively, these segmentation criteria shape how businesses formulate strategies, allocate resources, and respond to external challenges.
\\ \textit{Note:} correct because it identifies and elaborates on the three key segmentation criteria-organizational size, industry type, and geographical location-demonstrating how each one uniquely shapes a company's strategic approach and decision-making processes across various contexts.
\item \textbf{Answer:} John Cowan's family incurred a total meal cost of \$18.55, which likely includes the base price of the meals plus applicable sales tax. This total reflects the importance of understanding cost management within family budgets and personal finance. In business transactions, issuing separate checks can complicate accounting and payment processing, leading to inefficiencies and potential discrepancies in sales reporting. Conversely, a single check simplifies these processes, facilitating easier tracking of expenses and ensuring a cohesive payment method, which can enhance customer satisfaction and streamline operations. The choice between separate and single checks ultimately impacts both operational efficiency and customer experience in the hospitality sector.
\\ \textit{Note:} The answer correctly highlights the total meal cost's inclusion of sales tax while emphasizing the implications of payment methods on accounting efficiency and financial clarity in business transactions.
\end{enumerate}

\vfill
\noindent\textit{End of Answer Key}
\end{document}